\chapter[Introduction]
{An Introduction to System Administration}
\label{chap:introduction}

\begin{quote}
{\em system administrator, n.}:\\
one who, as a primary job function, manages computer and network systems
on behalf of another, such as an employer or client.\\
\end{quote}

\begin{figure}[hb]
	\raggedleft
	\includegraphics[width=.25\textwidth]{01/pics/we-are-happy-to-serve-you-coffee-mug}
	\caption[Anthora Coffee Cup]{System Administrators: well
		caffeinated and always happy to serve you.
		({\em previous page})
		\label{fig:we-are-happy-to-serve-you}}
\end{figure}

\section{What exactly does a System Administrator do?}
\label{introduction:introduction}

Every semester I start my class with this simple
question: ``What exactly does a System Administrator
do?''  My audience consists of undergradate and
graduate students of Computer Science or Computer
Engineering; some of them have experience working as a
System Administrator or IT support (either part-time
or full-time), while others may maintain their own
home network or otherwise perform some of the tasks
commonly associated with the job of a System
Administrator.  Some may have no practical experience
in this area, but all do have an idea of what the job
entails.

As a result, the responses are varied, ranging from
very specific (``Use \verb+tcpdump+ to troubleshoot
network problems.'') to the more general (``Maintain
the network.''), quickly illustrating that there is no
one concise description of the professional duties of
System Administrators; rather, we identify a number of
tasks that are expected to be performed by people
holding this title.  The broad range of these tasks --
especially when put in context of the size and goals
of different organizations and their deployment
footprint -- makes obvious the fact that a System
Administrator in one company may not do the same thing
as one in a different company.  In fact, even within a
single organization we usually find that two people
may well both be called ``SysAdmin'', but perform
rather different duties.

So what {\em does} a System Administrator do?  It
seems that most people agree that this job has to do
with computers in some capacity, but that really is
already where the definitions start to become fuzzy:
System Administrators are in charge of ``servers'' as
well as of personal computers, desktop machines,
laptops and, increasingly, mobile devices.  With the
evolution of computing coming full circle, we find
that the migration from the mainframe computer to the
personal and independent computing device to the
networked server back to the central {\em
Infrastructure as a Service}\index{Infrastructure!As a
Service} model inherent in today's Cloud
Computing\index{Cloud Computing} concepts places the
System Administrator in the middle of it all.

So perhaps the main focus of a System Administrator's
job is then really the central connection between the
independent systems, the network itself?  Hmmm, System
Administrators surely are involved in the deployment
and installation of the computers in a datacenter (or
is that the task of specific Datacenter Technicians?)
and connecting all the different components certainly
involves a lot of cables.  But then, don't we have
{\em Network} Administrators, or is that a more
specialized subcategory of System Administrators?

System Administrators seem to spend as much time
typing cryptic commands into dark terminal windows as
they do running cables and labelling hardware, and
while they may no longer shuffle punch cards, they
frequently do write programs to help them complete
their tasks.  System Administrators are known to get
woken up in the middle of the night when things go
``bump'', and as a result they are also known to have
a fondness for caffeinated drinks.  They are able to
kickstart a generator, assess the required cooling
power for their server room, use duct tape in creative
and unexpected ways, assemble servers out of
mysterious looking parts, and may end up handling a
circular saw or other heavy machinery when their
Leatherman multi-tool cannot complete the task.

System Administrators plan, budget and design networks
and backup or storage systems, add and delete users
(well, user {\em accounts}, anyway\footnote{There is a
strong love/hate relationship between System
Administrators and their users.  Much as the SA may
joke that they wish they could make users disappear,
without users the systems they are in charge of might
well hum along uninterrupted, but would ultimately be
entirely useless.}), install and update software
packages, draft policy documents, fight spam with one
hand while rebuilding a corrupted revision control
system with the other.  They have access to all
systems in the organization, may undergo retina- and
fingerprint scanners to access ``Mission
Impossible''-style protected datacenters and spend
countless hours in front of a multitude of computer
screens, typing away on oddly shaped keyboards
consuming not entirely healthy amounts of coffee and
energy drinks.

Well... in some places, a System Administrator might
do all of this.  In others, there might be different,
more specialized people for the various tasks:  there
might be datacenter technicians and ``SiteOps'',
Network Administrators, System Programmers, System
Architects, Operators, Service Engineers, Storage
Engineers, \glslink{sre}{Site Reliability
Engineers}\index{SRE}, Virtual Operations,
Infrastructure Architects... the number of different
job titles for things that might otherwise fall into
the more general ``SysAdmin'' category seems endless.

The various areas of expertise included in the day to
day routine such as system design, infrastructure
architecture, system fault diagnosis, hardware
benchmarking, and others eventually require experience
in a number of related fields; a lot of them involve
more than just a little bit of programming experience
and in some cases complex infrastructure tools based
on solid software engineering practices need to be
built.  This illustrates that just as it is hard to
clearly define where System Administration begins, we
often can't quite put our finger on where it ends and
another discipline becomes the primary job function.

Even if job descriptions and duties differ
significantly, there are a lot of people who simply
identify as {\em System Administrator} regardless of
job title.  We all have something in common: {\em we
manage and maintain computer and network systems}.

\begin{experience}
{\em A typical bug} \\

At one place of employment we had a tool that would
add users to a given host by gathering the login
information for all the users, \verb+ssh(1)+ to each
host and update \verb+/etc/passwd+ using the host's
native tools (such as \verb+useradd(8)+, etc.). Many a
system administrator has written a similar tool, and
for the most part, this program worked reasonably
well. (We will discuss how to better manage user
accounts across large numbers of machines later in
this book.) \\ [10pt]

But all software fails eventually, and for one reason
or another I had to debug the tool. Groveling through
a few hundred lines of perl that had clearly been
written by different people in different stages of
their career, I came across a code block that included
the following call: \\ [10pt]

\verb+chmod(0666, "/dev/null");+ \\ [10pt]

Asking around, it became clear that {\em somehow} some
of the systems ended up in a state where
{\tt /dev/null}\index{\tt /dev/null} was unreadable
or -writable by normal users, which leads to a number
of complications.  But nobody had been able to track
down just how exactly this happened until one day we
were lucky enough to witness the change in
permissions,  but the only activity by a user with
sufficient privileges to make such a change was a
\verb+sudo(8)+ invocation of \verb+less(1)+. \\
[10pt]

\verb+less(1)+ happens to make use of a history file
to remember search commands between invocations, much
like your shell may use a history file (such as
\verb+~/.bash_history+). \verb+less(1)+ is aware of
the fact that the history file might contain sensitive
information, and it actively changes the permissions
on that file to only allow read/write access to the
owner.  But for the root user, we explicitly symlinked
all common such history files to \verb+/dev/null+ to
avoid accidental leaking of secrets to disk. \\ [10pt]

And therein lies the bug: \verb+less(1)+ was invoked
by the super user, it would check the permissions on
the file \verb+/root/.lesshst+, follow the redirection
to \verb+/dev/null+ find that they're not \verb+0600+
and call \verb+chmod(2)+, yielding and
unreadable/unwritable \verb+/dev/null+.  We were able
to confirm this behaviour and identified a code change
in later versions of \verb+less(1)+ that fixed this
problem. \\ [10pt]

Understanding such failures and identifying their root
causes requires a deep understanding of the operating
system and its tools, but tracking down bugs like
this\cite{intro:devnull} is just one of the many
aspects of a System Administrator's daily routine.

\end{experience}

\section{The Profession of System Administration}
\label{introduction:profession}

Most professions have a fairly clear and formal
definition as well as an obvious career path and
background requirements.  In fact, for a person to
take on a particular line of work, it frequently is
a requirement to undergo a specific, regulated
education, to perform a certain amount of practical
training (such as an apprenticeship), to subscribe to
specific professional ethics\index{ethics} and
guidelines, and eventually to obtain a license to
practice the given profession.  Doctors, scientists
and lawyers all have well-defined professions, as do
engineers of most disciplines.\footnote{Software
Engineering, another young discipline related to
System Administration and technically a specialized
branch of Computer Engineering, differs in that it
does {\em not} have strict licensing regulations.
Anybody can call themselves a ``Software Engineer'',
just as anybody can call themselves a ``System
Administrator''.}

System Administration, on the other hand, does not
(yet) have any of the criteria given above: there is
no specific, regulated training, and as we have seen
in Section \ref{introduction:introduction} not even
agreement on what precisely the job definition ought
to be.  The \gls{lisa}\index{USENIX} attempts to fill
this void to some degree by providing for an informal
but professional community.  Still, due to the nature
of the changing environment, any attempt to create a
concise job description invariably fails by resorting
to very general terms.  Consider the following
description given by the Buerau of Labor
Statistics\cite{intro:bls}:

\begin{quote}
{\em Systems administrators} design, install, and
support an organization's computer systems. They are
responsible for LANs, WANs, network segments, and
Internet and intranet systems. They work in a variety
of environments, including large corporations, small
businesses, and government organizations. They install
and maintain network hardware and software, analyze
problems, and monitor networks to ensure their
availability to users. These workers gather data to
evaluate a system's performance, identify user needs,
and determine system and network requirements.
\end{quote}

While this captures {\em in general} the common job
duties of a System Administrator, it does not answer
the question we asked initially: Just what {\em
exactly} does a System Administrator do?

In order to better grasp the generic description, I've
found it useful to break down the title into its
components, {\em System} and {\em Administrator}.
What, then, is a ``System''?  Any dictionary will give
you a reasonable definition along the lines of:

\begin{quote}
{\em system, n}: A group of interacting, interrelated,
or interdependent elements that together form a
complex whole.
\end{quote}

In this context we will focus on {\em computer-human
systems}\cite{intro:burgess-analytical} consisting of
any number of computing devices, the network(s)
connecting them, the users utilizing these resources,
and of course the impact of the goals of the
organization running such systems.

What about ``Administrator''?  Merriam-Webster's
Collegiate Dictionary defines the term ``to
administer'' as:

\begin{quote}
{\em administer, v}: to manage or supervise the
execution, use, or conduct of
\end{quote}

The profession\index{System Administration!Profession}
of System Administration should therefore be described
as one in which practitioners ``manage or supervise
the execution, use, or conduct of a group of
interacting, interrelated, or interdependent
computer-human elements'', leaving us, in terms of
specificity, about where we began!  We find slightly
more formal definitions of actual job titles, desired
skills and a rather useful distinction between ``small
uniform'', ``complex'', and ``large and complex''
sites in the \gls{lisa} booklet ``Job Descriptions for
System Administrators''\cite{intro:lisa-jobs}, but in
searching for a suitably clear definition of the
profession, we realize that the formal education as a
System Administrator is almost completely lacking.

Job descriptions and postings normally list a degree
in Computer Science or Computer Engineering as
``desirable''; more frequently you will find a
requirement expressed as ``BS, MS or PhD in Computer
Science or equivalent work experience''.  Computer
Scientists will wonder just how, exactly, ``work
experience'' can be mapped, treated as equivalent to
their formal education -- all too often, ``Computer
Science'' in this context is equated with ``knows how
to program, understands TCP/IP''.

As invaluable as actual hands-on experience is, it is
no substitute for a formal degree in the area of
Computer Science.  (Nor is having such a degree a
substitute for practical experience.) Rare is the
candidate who can fully understand the average and
worst-case runtime of his or her programs and scripts
without having learned Big-O Notation; rare the system
programmer who completely grasps the elegance of
modularization provided by, say, function pointers
without having learned the $\lambda$-calculus.
Becoming a proficient System Administrator without
this background is certainly possible, but it makes it
significantly harder to advance beyond a certain
level.

Be that as it may, the focus on many years of
real-world experience as the primary qualification is
both cause and effect.  In the past, there really
existed {\em no} formal training whatever: technically
skilled people found themselves in the position of
being the go-to guy to fix the computers, to make
things work.  Perhaps they ended up working with a
more senior administrator in a sort of informal
apprenticeship, but for the most part their skills
were learned by experience, through trial by fire.
The lack of any formal program in this area
necessitates and then further encourages
self-learning, and many things simply can only really
be understood by experience.  It's a clich\'{e}, a
platitude that you learn the most from your worst
mistakes -- it is also entirely true.

As a result, many of the senior people in hiring
positions do not have an academic background and are
hesitant to require what was not available then to
fill current openings.  System Administrators are also
still frequently treated as -- and at times view
themselves as fulfilling -- a ``janitorial'' role.  If
no industry standard exists by which to measure a
candidate's accomplishments in the field, how can we
require formal education for a job that eludes
definition in the first place?

Therein, however, lies a fallacy.  Practical
experience and formal training do not make up for one
another; they complement each other.  It {\em is} true
that we learn the most from our worst mistakes (and
thus need a chance to make them); it is {\em also}
true that we frequently are only {\em able} to learn
due to a deeper understanding of the problem.  \\

The profession of System Administration is incredibly
interesting precisely {\em because} it eschews a fixed
definition, a formal governing body, a static career
path, or a licensing exam.  It attracts and invites
interesting people from all sorts of backgrounds; it
allows phenomenal potential for job growth and it is
-- must be, due to the advances in the industry --
fast-paced and continually evolving.  But it has
become more mature: professional organizations like
\gls{lisa} have contributed significantly to the
status of the profession by formalizing the
requirements and duties as well as providing a Code of
Ethics\footnote{For a more detailed discussion of both
the (lack of a) definition of the profession and the
ethical obligations we face, please refer to
\cite{intro:primum-non-nocere}.}, and over the last
few years more and more universities have started to
offer classes and degree programs in this area.

System Administrators may not need to be licensed by a
central board anytime soon -- and people will continue
to argue about whether or not that is desirable or
even possible in this field -- but what a few years
ago still was best described as ``a job'' certainly
has grown up to become a career, a craft, a {\em
profession}.

\section{System Administration Education}
\label{introduction:education}

System Administration is rarely taught as an academic
discipline in large part due to its perceived
``blue-collar'' status: from the top of the scholarly
ivory tower, it must be hard to distinguish where an
entry-level job like ``Tech Support'' or the ``Help
Desk'' ends and where a profession demanding a
significant background in Computer Science and
Engineering begins.

System Administrators are highly skilled information
technology specialists shouldering significant
responsibility in any organization.  As we discussed
in the previous sections, the variety of job
descriptions and differences in what a System
Administrator actually may be doing make it difficult
to provide simple step-by-step instructions on how to
enter this field.  As a result, System Administration
has long been a profession that is learned primarily
by experience, where people grow into a position in
order to fulfill the requirements of an organization
rather than follow a career path well-defined by
courses, degrees, and meaningful certifications.

The industry has responded by producing a large number
of practical certification exams that hope to attest
to the student's proficiency in the subject matter.
Practically, most of these certifications appear to
primarily test a person's ability to memorize specific
commands, further reinforcing the notion that System
Administration can be reduced to empirical knowledge,
making practical experience indeed equivalent, if not
superior to one of those scholastic degrees.

But certification that one remembers a specific
vendor's unique configuration file syntax does not
imply actual {\em learning} has taken place; holding a
one-week ``boot camp'' -- the name itself is revealing
of it's educational intention --  to drill enough
information into the participants' heads such that
they pass an exam at the end of the week does not
guarantee long-term retention of that knowledge.
Furthermore, there is no oversight of the topics
taught in such classes, no review of suitability or
correctness.

System Administrators who excel at their jobs do so
not because they have accumulated arcane tribal
knowledge about obscure pieces of hardware and odd
software systems (useful as that is), but because they
understand fundamental underlying principles and
combine that knowledge with their concentrated
experience.

System Administration -- its principles, fundamentals
and practice -- should be {\em taught}.  Since you are
reading this book, chances are that you are currently
taking a class in System Administration as a student
or perhaps you are teaching such a class.  But we
cannot teach an entire profession in a single class or
even using a few select courses.  It is necessary to
develop academic degree granting programs that
comprehensively prepare students for the varied and
wide requirements imposed on them when they enter the
industry.  Such programs should be combined with
extensive real-world and hands-on experience wherever
possible, perhaps by including internships,
cooperative education and apprenticeships provided by
industry leaders.  As students you are entitled to
pratical and useful exercises; as educators we have an
obligation to create these and make available the
required resources.

We need classes that combine the academic mission of
fostering independent research with the factual
requirements posed by the positions offered in the
marketplace.  We need instructors who have years of
experience in varied environments, who understand the
practical restrictions that may make the ideal
theoretical solution utopian but who are consciously
aware of these boundaries and able to teach their
students the awareness of them.  At the same time, we
need scholars willing to further this profession
through research and build a theoretical body of
knowledge to become the foundation of future programs.

\subsection{Challenges in System Administration Education}
\label{introduction:education:challenges}

Formal classes in System Administration as part of a
Computer Science or Engineering curriculum are still
uncommon, but in the past few years more and more
institutions have recognized the industry's need for
academic courses that adequately prepare students for
the multitude of responsibilities within this field
and have started to offer such classes (and in some
cases complete degree programs).  But the history of
this profession brings with it a number of unique
challenges when it comes to formally teaching its
principles and practices.

To really understand and appreciate some of the most
general aspects of System Administration, you need to
be exposed to actual running systems.  Practical
experience is so integral to this profession that it
cannot be separated from the theoretical knowledge and
postponed until the student enters his or her first
job or apprenticeship.  But therein lies a significant
hurdle to traditional teaching methods: students need
to administer a system, to have superuser access, to
have a chance to configure a system for a specific
service, and to make the kind of spectacular mistakes
that experienced System Administrators value (if only
in hindsight).

This normally conflicts with the requirements of the
IT department at your university: students would
require access to a number of different OS when the
school's system administrators strive for a certain
level of homogeneity. In order to understand OS
installation concepts, file system tuning, and other
low-level principles, students need to perform these
tasks themselves. Learning to debug network
connectivity issues or being able to actually see the
payload of captured network traffic requires access to
raw sockets, which the staff responsible for the
security of the campus network would certainly rather
not provide.

After years of struggling to give all students equal
access to the resources required for the diverse
assignments and frequently having to change the
assignments to be able to make use of available
resources, I finally decided to do what many companies
do: I outsourced the resource provision problem into
``the Cloud'' -- specifically, to Amazon's
\gls{ec2}\index{AWS!EC2}\footnote{In
this book, I will frequently make reference to or cite
examples from this experience.  At the same time, I
also will suggest exercises that cannot be done in a
cloud environment and that may require access to
dedicated hardware.}.\cite{intro:login-teaching-sa}
Similar results can be achieved by granting students
access to a dedicated laboratory environment or by
providing another on-demand infrastructure via
virtualization of the local compute resources.  If you
are a student, it is worth your time to investigate
all possible resources available to you; if you are an
instructor, consult with your IT department as to the
availability (or perhaps development) of such systems.
\\

But teaching System Administration should also combine
the practical elements noted above with topics from
other academic fields.  As an area of specialization,
System Administration fits in well within the current
Computer Science curricula, since the profession's
requirements draw heavily from many of the same
subjects.  Shared requirements include Operating
Systems, Networking, Database Systems, Distributed
Systems, Cryptography and of course Software
Engineering; at the same time, System Administrators
benefit from deep insights in Computer Architecture
and Systems Engineering as well as Project Management
and of course ``the 5 elements of administration''
(Planning, Organizing, Command (directing people),
Coordination, Control, also listed as
Budgeting)\cite{intro:wiki-fayolism}, usually covered
in Business Administration classes.

To advance the practice of System Administration, the
two existing professional organizations \gls{lisa} and
the \gls{lopsa} are already promoting best practices
and important communities.  Still, we are missing
strong proponents of System Administration as an
academic discipline\footnote{Mark Burgess, previously
of the Oslo University College in Norway, one of the
only universities offering a Master's degree in
Network and System Administration, deserves explicit
mention as one of the few exceptions.} But there are
more and more of us who believe that it's time for the
profession to be recognized as both requiring a
theoretical background in addition to extensive
practical experience and benefitting from the
opportunities of dedicated research opportunities and
the development of universally accepted degree
granting requirements.  Some even write books to
further this goal...

\section{The Three Pillars of Exceptional System Design}
\label{introduction:pillars}

More than just maintain computers and infrastructure
components, System Administrators control entire
systems.  With more experience and seniority, they are
ultimately in charge of building these systems, of
designing the infrastructure and its components so
that the job title frequently morphs into that of a
Systems or Network ``Architect''.  It seems that
people can't help but draw parallels to the so-called
``real'' world, where engineers and architects have
well defined roles (and status).  This does not come
as a surprise; as Brooks\cite{intro:brooks-design}
illustrates masterfully, the parallels between
designing and building a house, a computer or a
complex piece of software are plentiful.

Billy Hollis\index[names]{Hollis, Billy} points
out\cite{intro:hollis} that the title of the
``Software Architect'' is a misleading analogy, that
the duties traditionally performed by people holding
this title are much more akin to those of a structural
engineer than those of a (traditional) architect.  The
role of the System Administrator is likewise
frequently comparable to that of the structural
engineer in the mechanical world, but the advanced
senior System Administrator -- who ultimately steps
into the role of planning, designing and overseeing
the construction of complex systems -- might better
known as ``Systems Architect''.  Structural systems
engineers are then needed to aide in the final design
and implementation of the system in question.

Conceiving of and designing systems as complex as
those System Administrators are in charge of, with
their myriad pieces, the unpredictable human element
and the rather unfortunate tendency of computers to do
precisely what they were instructed to do rather than
what we intended them to do requires expert in-depth
knowledge in a broad range of topics.  Exceptional
System Administrators are generalists -- they have
been trained in a wide area of topics, deepened their
understanding of several and perhaps have become
subject matter experts in a few, but ultimately back
into a role that allows them to apply their ability to
connect the dots, to see the big picture and the
connections between and effects of multiple systems on
each other.  They are well-equipped to model new
components of an infrastructure or to redesign a
system from scratch to meet the new requirements.

For a system to be suitable both for the immediate
requirements today and for a few years down the road,
when the current needs have changed, it must embody
two basic principles: Scalability and Security.
Neither of these can be added to a system after it has
been built:  trying to apply ``security'' after the
system interfaces have been defined yields
restrictions, limitations; trying to make a system
with inherent limitations perform under circumstances
it was not designed for yields hacks and workarounds
-- the end result frequently resembles a fragile house
of cards more than a solid reliable structure.

The third fundamental feature of expertly designed
systems, Simplicity, is simultaneously obvious and
counter-intuitive.  Simplicity underlies both
scalability and security, since reduced complexity
implies better defined interfaces, minimized ambiguity
in communications or data handling and increased
flexibility.  As with the other two core aspects,
simplicity cannot be added after the fact; it must be
inherent in the architecture.  Simplicity is the
enabler of both scalability and security.

An exceptional system exhibits inherent structural
integrity (another wonderful analogy to the world of
bricks and mortar), and this integrity is provided by
these three pillars.  Our focus on these components
may initially seem arbitrary: meeting requirements
across different teams with different priorities has
long been the bane of many a program manager due to
each team focusing on different qualities, of which we
could have picked any number.  However, upon more
detailed analysis and with some years of experience we
have time and again found Simplicity, Security and
Scalability to be the core qualities enabling a
harmonious development- and deployment process.

Throughout this book, we will analyze the systems we
discuss with special focus on these crucial criteria.
You will find that we will frequently talk in broad
abstracts or surprising detail about the implications
of our design decision.

\subsection{Scalability}
\label{introduction:pillars:scalability}

In recent years, the word ``scalability'' has become
one of the defining requirements for virtually all
technical solutions.  Providers of infrastructure
services and products throw it around in an attempt to
impress their customers by how much load their systems
can handle.  It typically seems near synonymous with
``high performance'', the ability to handle large
amounts of data, traffic, connections, and the like.
This may be misleading, as people might be inclined to
believe that such capabilities come at no additional
(or incremental or proportional) cost.  In fact,
scalability does not relate to the costs of the
running system, but to its architecture.  A scalable
system may well require additional resources and
throwing money at a problem actually {\em is} a
perfectly viable solution in many cases.  However, a
scalable architecture is one that does not require a
refactoring of the whole system, a restructuring of
{\em how} data is handled based on changing
circumstances.  Instead, it abstracts data from
process flow and simply adapts.

System Administrators frequently suffer from
Goldilocks Syndrome, striving to provide a solution
that is {\em just right}; wasting resources is just as
anathema to their definition of a correct solution as
not providing adaequate capabilities.  Therefore, we
tend to focus not only on the ability to scale {\em
up} -- to accommodate larger amounts of data or more
requests per minute, for example -- but also to scale
{\em down}, to free up unused resources based on the
demands.  As such, our definition of ``scalability''
leans more towards the overall flexible nature of a
scalable system, its ability to adapt to changing
requirements at run time; a term frequently associated
with Cloud Computing, ``elasticity'', perhaps more
aptly describes this feature, yet it fails to quite
capture as well the sense of meeting extreme
requirements.  We shall continue to use the term
``scalability'' for that reason.  \\

While many systems are {\em flexible}, few are able to
handle input or requirements an order of magnitude
different from those initially conceived of when the
system was designed.  In order to accomodate such
changes, systems are said to either scale {\em
vertically} -- that is, one system is able to handle
the added load, possibly by addition of certain
resources (network bandwidth, CPU power, memory, disk
space, ...) -- or to scale {\em horizontally} -- that
is, a single system is replicated and demands are
spread across them evenly.  Either approach has
certain implications on the overall architecture.  For
horizontal scaling to be beneficial, for example, the
problem space needs to be such that a distribution is
both possible and algorithmically feasible.  However,
this adds communication overhead, as interfaces and
data flow becomes more complex.

Whether by vertical or horizontal means, a scalable
system is one that readily adapts to changing demand.
The designer may not know what will be required of it
in the future, but will make choices that permit the
system to grow and shrink without bumping into
arbitrary limits.

We will take a closer look at the implications of this
idea on the software development practices and overall
Unix Philosophy in Chapter \ref{chap:unix}; for the
time being, let us consider the Unix tradition of
simple tools operating on streams of text as a
wonderful example of a clear interface definition.
Anybody developing a tool that accepts input only from
a file restricts the flexibility of the tool; this
frequently goes hand in hand with an implicit
limitation on the amount of data that can be handled
(think maximum file size, buffers frequently allocated
to read in the entire file, the ability to
\verb+seek(2)+ on the file handle, ...).  By choosing
text rather than a binary format, any future use of
the output is not limited by the original author's
imagination of what future users might wish to
accomplish.  Given sufficiently large or, rather,
diverse datasets, building more complex systems that
perform equally well under heavy load using
complicated, yet limited interfaces so frequently
developed by major software companies easily becomes a
frustrating exercise in determining these
boundaries.

But scalability is not only of concern when tools are
developed; as we design infrastructure components, we
need to be acutely aware of the interfaces between
them and what kind of data can flow back and forth.
Frequently we have no control over the input (both
type and amount), so our systems need to be fault
tolerant of many unenvisioned circumstances.  Even
though somewhat counter-intuitive at first, I argue
throughout this book that a robust system will remain
resilient to overall failure by being comprised of
infrastructure components that themselves may fail
quickly (and explicitly) and that such well-defined
behaviour underlies true scalability.

Hindsight being 20/20, scalability related issues
often can be traced back to a lack of imagination or
confidence in the system.  What if our initial use
case increases not by a factor of two or three, but
hundredfold?  Suppose our system is still in use in
five years -- will average input be likely to remain
the same?  In our discussions, exercises and problems
we will encourage students to consider how a system
performs if circumstances and inputs change by an
order of magnitude.  The ability to anticipate
extraordinary change in requirements requires some
practice and experience by the System Administrators
in charge of the infrastructure; reacting in
well-defined ways to such change is one of the core
features of a reliable system, and the principles that
make a system scalable -- fault tolerance, data
abstraction, clearly defined interfaces -- must be
applied at all stages.  We will reinforce these axioms
throughout all chapters in this book.


\subsection{Security}
\label{introduction:pillars:security}

All too frequently, the software or information
technology industry treats system security as an
afterthought, as something that can be added to the
final product once it has met all the other functional
requirements, after the user interface has been
determined and after all code has been written.  It is
then not surprising that the old adage that security
and usability are directly and inversely related seems
to hold true.  Nevertheless, I would argue that the
very problem statement -- ``the more secure you make
something, the less usable it becomes''\footnote{We
will elaborate on the somewhat surprisingly accurate
corrolary that ``The more secure you make something,
the less secure it
becomes.''\cite{intro:norman-security} detail in a
later chapter.} -- reflects a fundamental
misunderstanding, as it implies that usability is
present first and security ``added'' afterwards.

This approach suggests that the only way to reduce
risks is to {\em take away} functionality; but any
time you do that or otherwise restrict the users, they
will either stop using the product/system altogether
or come up with a creative solution that works around
your newly imposed restrictions.  To be effective,
security needs to be built into the system from the
design phase on.  That is, rather than starting out
with a solution that provides the desired
functionality and then attempting to figure out how to
get it to a secure state, we should instead begin with
a secure albeit restricted state and then slowly add
functionality -- without compromising safety -- until
the desired capabilities are available.  That is, we
need to view security as an enabling factor present at
the design's conception.

Much like software, a system infrastructure is usually
developed with much looser an idea of what one wishes
to accomplish than people are willing to admit.  In
order to achieve maximum functionality and the widest
possible use, people tend to design interfaces with
few restrictions.  Often it is the role of the
security-minded engineers to ask questions that
require the developers or designers to revisit their
requirements; proactive security imposes restrictions
on what the system is able to do.  That is a {\em
good} thing!  General-purpose systems are much harder
to design than special-purpose sytems (``Constraints
are friends''\cite{intro:brooks-design} whitelists
provide significantly better and more reliable
security than blacklists.

Any parent can tell you that it is nearly impossible
to take away a toy from a child, even if the toy has
not been used in months or years.  The very idea of
something being taken away seems to trigger in people
(of all ages) a ``but I might need it some day''
response.  Giving users access to a resource (network
access, disk space, physical access, CPU power, ...)
is trivial -- restricting the use once access has been
granted is near impossible; the genie cannot be put
back into the bottle.  It is therefore imperative to
understand precisely what level of access to what
resources you {\em actually} need and not build a
system for the most widest possible use.  Well defined
-- restricted -- small components can still provide
the flexibility to build scalable systems, but each
component needs to be designed from the beginning with
security in mind.  Just as a chain is only as strong
as its weakest link, an infrastructure is only as
secure as its most open or exposed component.  \\

Rare is the software developer who dismisses
``security'' as of minor importance, yet this mythical
property everybody pays at least lip service to has an
incredibly wide definition and may mean different
things in different contexts.  Generally speaking, we
require {\em risk management} (we discuss this in more
detail in a later chapter); more specifically we will
look at a number of computer-human system specific
aspects of security such as (but not limited to):

\begin{itemize}
	\item Cryptography\index{cryptography}, and the three main features that allow
		it to help mitigate risks:
		\begin{itemize}
			\item Secrecy or Confidentiality
			\item Accuracy or Integrity
			\item Authenticity
		\end{itemize}
	\item Physical Security
	\item Service Availability
	\item Service Design
	\item Social Engineering and other aspects of human nature
	\item Trust
\end{itemize}

Throughout this book, we will allude to these
components wherever appropriate, and we encourage
students to come up with possible threat scenarios as
a regular exercise.  Establishing and fostering a
proactively defensive, some may say ``paranoid'',
mindset is an important element of a System
Administrator's development.  Knowing that we cannot
provide an all-encompassing mantle of security after
the system has been designed (much less deployed!), we
strive to integrate it in all stages, leading it to
bolster our infrastructure as an inherent feature.



\subsection{Simplicity}
\label{introduction:pillars:simplicity}

As our discussion of scalability and security
suggested, the practical application of these
principles yields a reduction of interfaces,
end-points, use cases and overall variance.  In other
words, scalable and secure systems are less {\em
complex} and -- it is worth drawing this distinction
explicitly -- much less {\em complicated}.  A complex
system may be well-organized and exhibit a clear,
logical structure yet require subtle or intricate
connections or components.  Complicated systems, on
the other hand, are irregular, unpredictable or
difficult to follow.

``Complexity is the enemy.'' This quote has been
repeated endlessly in so many contexts that it can
almost be assumed common knowledge amongst software
engineers, cryptographers and other security
experts\cite{intro:schneier-ce}.  But while many
people in the information technology industry may
agree on this in principle, I have found that, like so
many aphorisms, the translation into the real world --
the actual application of its consequences -- trails
far behind the good intentions.

In the world of System Administration reducing
complexity is particularly difficult.  As we discussed
in Sections \ref{introduction:introduction} and
\ref{introduction:profession}, managing large
computer networks has {\em inherent} complexity:
multiple components must interface with each other in
many ways to help people -- simultaneously the origin
of true entropy as far as computers are concerned and
dangerously predictable -- accomplish their individual
and collective goals.  But it is important to
differentiate between {\em required}
complexity\index{Complexity!essential} and {\em
accidental}
complexity\index{Complexity!accidental}\cite{intro:brooks-mmm}.
We strive to reduce overall complexity by building
ultimately intricate systems out of smaller, simpler
components.

To repeat an almost ubiquitous prime example of how
simplicity enables astounding flexibility and may
allow you to build complex systems, we will time and
again draw the analogy to the concept of toy building
blocks: themselves simple, nearly unbreakable and
available in a reasonably limited number of shapes and
sizes you can build just about anything with them.

Like these little blocks, the Unix operating system
builds on the philosophy of simple tools that ``do one
thing and do it well'', that work together, commonly
by being connected via the ubiquitous pipe,
operating on text
streams\cite{intro:mcilroy-unix-philosophy}.  We aim
for solutions that exhibit comparable elegance.
``Simple'' does not, it should be noted, mean
``trivial'' or ``easy'': it is so much easier to {\em
add} features, to increase the output, to justify
additional input than it is to reduce a component to
its bare necessities.  As Antoine de Saint
Exup\'{e}ry\index[names]{de Saint
Exup\'{e}ry, Antoine},
observed\cite{intro:exupery:terre-des-hommes}:

\begin{quote}
{\em Perfection is reached not when there's nothing left to
add, but when there's nothing left to remove.}
\end{quote}

But how do we design a simple infrastructure?  The
simple answer is: we can't.  The end result will
always be complex.  But the systems we create will be
more fault tolerant, more performant, more flexible
-- in short: better -- if we remember that our
building blocks need to resemble Lego blocks more than
swiss army knifes.  In other words: we design {\em
for} simplicity.  We keep interfaces internally
consistent and logical.  We use simplicity as a
scalpel to cut away the clutter.

The ability to resist the temptation to ``enhance''
your architecture or software tool to implement
additional features that, if at all necessary, might
better be provided by a separate component, usually
only comes with strong will power and years of
experience.  At the same time, simplicity is found to
underly all of the exemplary technologies you are
likely to encounter and some of which we will present
in the coming chapters.  Sometimes it takes a second
look to recognize and appreciate the genius of
simplicity; hopefully this book will help sharpen your
awareness of this trait.

%"Fools ignore complexity. Pragmatists suffer it. Some can avoid it.
%Geniuses remove it." -- Alan Perlis

\section{The Future of System Administration}
\label{introduction:future}

The System Adminstrator's job definition is diverse,
and so is the rather significant change the profession
has undergone over the years.  Gone are the days of
shuffling punch cards, of cycling magentic tapes, of
connecting your dumb terminals using RS-232 serial
connections, of stringing ribbon cables inside large
computer cases... or are they?

Some things haven't really changed all that much.
Technologies come and go, but certain principles have
remained the same.  Using fibre-optic cables rather
than twisted-pair, using Infiniband or Fibre Channel
instead of Ethernet or parallel \glslink{scsi}{SCSI}, or
using wireless networks instead of physically
connected systems does not fundamentally change the
day-to-day operations.

Virtualization\index{Virtualization} and Cloud
Computing\index{Cloud Computing} may seem to limit the
need of an organization to hire their own System
Administrators, but I believe that rather than a
threat to the profession these technologies are simply
one of the ways in which we make progress as
information technology specialists.  A chapter towards
the end of the book discusses the future direction of
System Administration and revisits these topics as
well as some of the latest trends in how System
Administration intertwines with other disciplines --
the terms ``DevOps\index{DevOps}'' and ``Agile
Infrastructure'' (or ``Agile Operations'') deserve
particular attention in that context.

A few years ago, the Personal Computer signified a
paradigm change away from the central Mainframe
computer accessed remotely by users towards a
distributed storage and compute model.  System
Administrators found themselves supporting more and
more individually powerful connected devices, trying
to centralize certain resources and services such as
shared filespace or regular backups, for example.  In
the recent past, the move towards software and data
storage services ``in the cloud'' has brought the
evolution full circle:  once again, we are challenged
to provide reliable central services that are accessed
by small, though increasingly mobile, devices.

The separation of {\em systems} from {\em networks},
if it ever really existed, is disappearing.
Standalone systems separated from the Internet are
rare; systems operating completely on their own
without {\em any} network connectivity have become
largely impractical.  Does this change the job of the
System Administrator?  Does it make things more or
less interesting, difficult, challenging, exciting,
different?  Will we see another shift back towards a
model where data storage and processing happens again
on the edge nodes of the network?

We don't know.  But we do know this: whatever happens,
System Administrators will be needed.  Perhaps they
will carry a different title and perform different
practical tasks, but the concept of the profession
will remain.  Organizations small and large will need
somebody with the knowledge and experience to analyze,
troubleshoot and design new infrastructures; somebody
who builds and maintains scalable and secure systems;
somebody with an appreciation of the simple things in
life.

\begin{quote}
{\em Plus \c{c}a change, plus c'est la m\^{e}me chose.}
-- Jean-Baptiste Alphonse Karr\index[names]{Karr, Jean-Baptiste Alphonse}
\end{quote}
\vfill

\begin{figure}[hb]
	\includegraphics[width=1\textwidth]{01/pics/Ibm704}
	\caption[An IBM 704 mainframe]{An IBM 704 mainframe.
		With cloud computing coming full
		circle, we already live in the future!
		\label{fig:mainframe}}
\end{figure}

\pagebreak

\chapter*{Problems and Exercises}
\addcontentsline{toc}{chapter}{Problems and Exercises}
\section*{Problems}
\addcontentsline{toc}{section}{Problems}

\begin{enumerate}

\item Create a course notebook (electronic or
otherwise).  In it, write down your notes about each
chapter, add any links to additional information,
noteworthy insights, etc.  After each chapter, write
down lessons learned.  Differentiate between those
directly applicable to you and those you consider
worthwhile reviewing or investigating in the future.

\item Create a folder for your course work as a
central place for all your documentation.  Whenever
you go through practical exercises, make sure to write
down how you solved a specific problem.  Your
documentation will later on help you review the
lessons you learned and can be used as a practical
how-to guide, so make sure to verify all steps before
you write them down.

\item Ask a System Administrator (for example: in your
university, of an open source project you participate
in, at work) to describe their daily routine.  Compare
notes with others.

\item Consider the current most popular websites and
Internet businesses or organizations.  Try to find out
details about their infrastructure and how it is
maintained.

\item Research available classes and degree granting
programs in the field of System Administration.
Compare their contents to those of industry standard
certification courses.

\item Research professional organizations related to
System Administration (such as LISA and
\glslink{lopsa}{LOPSA}).

\begin{enumerate}

\item Review their mission and content and consider
joining their mailing lists.  They provide a great way
to keep up to date with real-world experiences and
itemlems.

\item How do these organizations compare to the
\glslink{acm}{ACM}, \glslink{ieee}{IEEE} or
\glslink{isoc}{ISOC}?

\end{enumerate}

\item Research common Internet standard bodies and
practices.  What is an RFC?  What do the
\glslink{ietf}{IETF}, \glslink{iana}{IANA} and
\glslink{icann}{ICANN} do?

\item Consider the systems you have access to: what
are their primary functions, for what goals were they
designed?  Suppose they grew by an order of magnitude
-- what itemlems would you foresee?

\item Consider the systems you have access to: what
kind and how is access granted?  What kind of security
itemlems can you imagine?

\item Research the terms ``DevOps''\index{DevOps} and
``SRE''\glslink{sre}{SRE}.  In how far do the practices it
represents change the typical job description of a
System Administrator?

\end{enumerate}

\section*{Exercises}
\addcontentsline{toc}{section}{Exercises}

\begin{enumerate}

\item
\label{item:basic}
Practice basic system tasks in your virtualized
environment by creating different OS instances and run
them.  Once running, log in on the virtual host and
run the required commands to:

\begin{enumerate}

\item display basic system information (\verb+uname+,
\verb+ifconfig+, \verb+netstat+, ...)

\item
display the partition table

\item
display the currently mounted file systems

\item
display available disk space
\end{enumerate}

\item Repeat Exercise \ref{item:basic} for a different
Unix flavor than what you are used to.  For example,
if you used a Linux instance for Exercise
\ref{item:basic}, repeat it using OpenSolaris or
FreeBSD.

\item Determine the Unix commands you execute most
frequently (for example via analysis of your shell's
history file).  Analyze the top three commands for
their complexity and interfaces.

\item Analyze your interactive shell usage.
\begin{enumerate}

\item Change your login shell for at least a week to
one that you are not accustomed to.  Make note of the
changes in behaviour.

\item Disable tab completion and any aliases you have
defined in your login shell for at least a week.  How
does this change your workflow?

\item Review the detailed documentation of your login
shell's \verb+builtin+s.  Practice the use of those
that you are not familiar with.

\end{enumerate}

\end{enumerate}

\pagebreak

\bibliographystyle{plainnat}
\begin{thebibliography}{99}

\bibitem{intro:devnull}``less bug, more /dev/null'',
Jan Schaumann, on the Internet at
\url{http://is.gd/NpihsH} (visited November 29,
2015)

\bibitem{intro:bls}Bureau of Labor Statistics, U.S.  Department of Labor,
{\em Occupational Outlook Handbook, 2010-11 Edition, Computer Network, Systems,
and Database Administrators}, on the Internet at
\url{http://www.bls.gov/oco/ocos305.htm} (visited December 10, 2011)

\bibitem{intro:burgess-analytical}Mark Burgess, {\em Analytical Network and
System Administration: Managing Human-Computer Systems}, Wiley \& Sons,
2004

\bibitem{intro:m-w}{\em Merriam-Webster's Collegiate Dictionary}, 10th Edition,
Merriam-Webster

\bibitem{intro:lisa-jobs}Tina Darmohray, {\em Short Topics in System
Administration: Job Descriptions for System Administrators}, 3d Edition,
USENIX Association, Berkeley, CA, 2010

\bibitem{intro:login-teaching-sa}Jan Schaumann, ``Teaching System
Administration in the Cloud'' in {\em ;login: The USENIX Magazine},
October 2010, Volume 35, Number 5

\bibitem{intro:wiki-fayolism} {\em Fayol's Elements of Management} on the
Internet at \url{https://en.wikipedia.org/wiki/Fayolism} (visited December
17, 2011)

\bibitem{intro:brooks-design}Frederick P. Brooks, Jr., {\em The Design of
Design}, Addison-Wesley Professional, 2010

\bibitem{intro:hollis} {\em Billy Hollis Still Builds
Apps}, transcript of an interview with Billy Hollis,
``.NET Rocks'', on the Internet at
\url{http://is.gd/Suk2hr} (visited December 23, 2011)

\bibitem{intro:norman-security}Donald A. Norman, ``THE WAY I SEE IT: When
security gets in the way'', in {\em Interactions}, November 2009, Volume
16, Issue 6, ACM, New York, NY

\bibitem{intro:schneier-ce}Niels Ferguson, Bruce Schneier, and Tadayoshi Kohno,
{\em Cryptography Engineering}, John Wiley \& Sons, March 2010

\bibitem{intro:brooks-mmm}Frederick P. Brooks, Kr., {\em The Mythical
Man-Month: Essays on Software Engineering, Anniversary Edition},
Addison-Wesley Professional, Anniversary edition (August 12, 1995)

\bibitem{intro:mcilroy-unix-philosophy}M. D. McIlroy, E. N. Pinson, and B. A.
Tague {\em Unix Time-Sharing System Forward}, The Bell System Technical
Journal. Bell Laboratories, 1978

\bibitem{intro:exupery:terre-des-hommes}Antoine de
Saint-Exup\'{e}ry, {\em Wind, Sand and Stars},
Harcourt, 2002

\bibitem{intro:primum-non-nocere}Jan Schaumann, {\em
Primum non nocere - Ethical Obligations in Internet
Operations}, VelocityConf NYC 2015; on the Internet at
\url{https://www.netmeister.org/blog/primum-non-nocere.html}
(visited January 16th, 2017)

\end{thebibliography}
